\clearpage
\section{结果分析}
\subsection{Depth Anything V2：不同输入尺寸的结果与分析}

为比较输入尺寸的影响，我们使用同一个 Depth Anything V2 模型和相同的 8 张图像，分别以 518 和 280 进行预测。两组结果采用相同的深度对齐与评估方法，并统一深度图的颜色范围，比较图像效果、误差和耗时。

\subsubsection{深度预测对比}

\ReportFigure[548bp]{relative_depth_grid.png}{8 张图像的深度预测对比。每组三行依次为原图、518 和 280 输入尺寸的预测；深度经对齐后统一着色，颜色越深表示距离越远。}

两种输入尺寸都能区分墙面、地面和家具的前后关系，整体结构较接近。518 在部分细节上更清楚，例如 0401 的婴儿床栏杆、0801 的椅背和桌椅边缘；280 的这些区域更平滑。

\clearpage

\subsubsection{误差与耗时比较}

\ReportTableCaption{两种输入尺寸的逐图误差及均值（对齐后；粗体为同一行、同一指标的较小值）}
\begin{ReportTable}{lCCCCCC}
\rowcolor{ReportNavy}
\ReportHead{图像编号} & \multicolumn{2}{c}{\ReportHead{AbsRel}} & \multicolumn{2}{c}{\ReportHead{深度 RMSE / m}} & \multicolumn{2}{c}{\ReportHead{XYZ RMSE / m}} \\
\rowcolor{ReportNavy}
& \ReportHead{518} & \ReportHead{280} & \ReportHead{518} & \ReportHead{280} & \ReportHead{518} & \ReportHead{280} \\
0001 & 0.0262 & \textbf{0.0210} & 0.0974 & \textbf{0.0830} & 0.1067 & \textbf{0.0906} \\
0101 & \textbf{0.1657} & 0.1823 & 0.6835 & \textbf{0.6230} & 0.7281 & \textbf{0.6647} \\
0201 & \textbf{0.1026} & 0.1588 & \textbf{0.8290} & 1.2179 & \textbf{0.8749} & 1.3046 \\
0401 & 0.0559 & \textbf{0.0394} & 0.1821 & \textbf{0.1511} & 0.1958 & \textbf{0.1627} \\
0601 & \textbf{0.0381} & 0.0386 & \textbf{0.2225} & 0.2445 & \textbf{0.2368} & 0.2622 \\
0801 & 0.0943 & \textbf{0.0749} & 0.5776 & \textbf{0.4644} & 0.6177 & \textbf{0.4908} \\
1001 & 0.0699 & \textbf{0.0640} & 0.2238 & \textbf{0.2195} & 0.2422 & \textbf{0.2359} \\
1201 & \textbf{0.0513} & 0.0638 & \textbf{0.1744} & 0.2123 & \textbf{0.1865} & 0.2283 \\
\rowcolor{ReportSummary}\textbf{均值} & \textbf{0.0755} & 0.0804 & \textbf{0.3738} & 0.4020 & \textbf{0.3986} & 0.4300 \\
\end{ReportTable}


\textbf{误差比较。}我们对每张图分别计算三项误差，再取 8 张图的平均值。518 的平均 AbsRel、深度 RMSE 和 XYZ RMSE 相比 280 分别降低约 6.0\%、7.0\% 和 7.3\%，整体上略有改善。

不过，\textbf{输入尺寸更大并不意味着每张图都更准确}。例如 0201 的深度 RMSE 从 1.2179 m 降至 0.8290 m，改善较明显；0001 和 0401 的三项误差则都是 280 更低。因此，需要结合多张图的结果判断，不能只看某一张图的效果。

\ReportTableCaption{两种输入尺寸的推理耗时（粗体为较短耗时）}
\begin{ReportTable}{CCCC}
\rowcolor{ReportNavy}
\ReportHead{输入尺寸} & \ReportHead{实际输入宽高} & \ReportHead{平均耗时 / ms} & \ReportHead{8 张总耗时 / s} \\
518 & $686\times518$ & 140.7 & 1.126 \\
280 & $378\times280$ & \textbf{53.3} & \textbf{0.426} \\
\end{ReportTable}


\textbf{耗时比较。}本次运行中，518 的平均耗时约为 280 的 2.64 倍。518 的平均误差略低、部分细节更清楚，而 280 用时更短，且仍能保留主要场景结构；两者体现了精度与速度之间的取舍。

\subsubsection{点云重建示意}

\ReportFigure[166bp]{relative_pointcloud_pair.png}{0001 的点云示意：左为参考深度生成的点云，右为 518 输入尺寸的预测点云。两者使用相同视角与显示比例，并用原图颜色着色。}

预测点云能够重建墙面、座椅和柜台的基本位置关系，与参考点云的整体形状接近，但局部轮廓仍有偏差。点云图主要用于展示重建效果，尺寸变化带来的精度差异以表中的 XYZ RMSE 为准。
